\section{Conclusion}

We presented Open-AoE, an open data infrastructure that couples approximately 2,000 hours of real-world egocentric manipulation data captured with consumer smartphones with a complete path from data to models.
On the production side, the capture and processing pipeline turns raw smartphone video into structured samples aligned across vision, language, hand motion, camera trajectories, and action boundaries; on the consumption side, the open-source toolchain supports synchronized visualization, 4D reconstruction, cross-embodiment retargeting, and training representations for Vision-Language-Action policies, World Action Models, and World Models.
Its broad coverage of scenes, actions, interaction objects, participants, device models, and fields of view allows this infrastructure to connect low-cost real-world collection with multiple embodied-learning paradigms.

Open-AoE is intended not as a static data package, but as an open ecosystem co-developed by data contributors, universities, research institutions, robot developers, and model teams.
We invite the community to contribute new real-world scenes and tasks, reconstruction and retargeting methods, training adapters, and evaluation results, allowing data, tools, and models to evolve together in an open loop.
Our goal is to make Open-AoE the lowest-barrier egocentric data infrastructure for embodied intelligence---making egocentric data easy for anyone to collect and use, so that data is no longer the bottleneck for embodied foundation models.
